\documentclass[aspectratio=169,10pt]{beamer}

% Shared Beamer setup for Fall 2026 course slides.
% Clean Cal Poly-inspired theme: no custom class files or uncommon packages.
\mode<presentation>{
  \usetheme{default}
  \usefonttheme{professionalfonts}
}

\definecolor{CalPolyGreen}{HTML}{154734}
\definecolor{CalPolyGold}{HTML}{BD8B13}
\definecolor{CalPolySage}{HTML}{7FA28F}
\definecolor{CalPolyMint}{HTML}{DDF1E7}
\definecolor{CalPolySand}{HTML}{A29F86}
\definecolor{CalPolyBeige}{HTML}{EFEEDE}
\definecolor{DeckInk}{HTML}{1F2933}
\definecolor{DeckMuted}{HTML}{5F6B7A}
\definecolor{DeckSoft}{HTML}{F7F7F5}

\setbeamersize{text margin left=0.55in,text margin right=0.55in}
\setbeamercolor{normal text}{fg=DeckInk,bg=white}
\setbeamercolor{structure}{fg=CalPolyGreen}
\setbeamercolor{title}{fg=white}
\setbeamercolor{frametitle}{fg=CalPolyGreen,bg=white}
\setbeamercolor{block title}{bg=CalPolySage,fg=white}
\setbeamercolor{block body}{bg=CalPolyMint,fg=black}
\setbeamercolor{alerted text}{fg=CalPolyGold}

\setbeamerfont{title}{size=\Huge,series=\bfseries}
\setbeamerfont{frametitle}{size=\Large,series=\bfseries}
\setbeamerfont{block title}{size=\normalsize,series=\bfseries}

\setbeamertemplate{navigation symbols}{}
\setbeamertemplate{itemize item}{\color{CalPolyGold}$\blacktriangleright$}
\setbeamertemplate{itemize subitem}{\color{CalPolyGreen}$\bullet$}
\setbeamertemplate{footline}{
  \hfill{\color{DeckMuted}\scriptsize\insertframenumber/\inserttotalframenumber}
  \hspace{0.18in}\vspace{0.12in}
}
\setbeamertemplate{frametitle}{
  \vspace{0.18in}
  \hspace{0.02in}\insertframetitle\par
  \vspace{0.08in}
  \noindent\makebox[\linewidth][l]{%
    {\color{CalPolyGold}\rule{0.55in}{2pt}}%
    {\color{CalPolyGreen}\rule{\dimexpr\linewidth-0.55in\relax}{2pt}}%
  }
}

\newcommand{\CourseNumber}{Course}
\newcommand{\CourseTitle}{Course Title}
\newcommand{\LectureNumber}{0}
\newcommand{\LectureTitle}{Lecture Title}
\newcommand{\LectureDate}{Date}
\newcommand{\CourseUnit}{Unit}

\newcommand{\coursenumber}[1]{\renewcommand{\CourseNumber}{#1}}
\newcommand{\coursetitle}[1]{\renewcommand{\CourseTitle}{#1}}
\newcommand{\lecturenumber}[1]{\renewcommand{\LectureNumber}{#1}}
\newcommand{\lecturetitle}[1]{\renewcommand{\LectureTitle}{#1}}
\newcommand{\lecturedate}[1]{\renewcommand{\LectureDate}{#1}}
\newcommand{\courseunit}[1]{\renewcommand{\CourseUnit}{#1}}

\author{Kirk Duran}
\institute{California Polytechnic State University, San Luis Obispo}
\date{\LectureDate}

\newcommand{\makedecktitle}{
  \begingroup
  \setbeamercolor{background canvas}{bg=CalPolyGreen}
  \begin{frame}[plain]
    \vfill
    \begin{center}
      {\usebeamerfont{title}\color{white}\LectureTitle\par}
      \vspace{0.18in}
      {\Large\color{white}Lecture \LectureNumber\par}
    \end{center}
    \vfill
    \noindent{\color{white}\rule{\linewidth}{1.2pt}}\par
    \vspace{0.08in}
    \noindent\colorbox{CalPolyGold}{\makebox[0.24\linewidth][c]{\color{white}\Large\CourseNumber}}
    \hspace{0.12in}{\color{white}\Large Kirk Duran}
  \end{frame}
  \endgroup
}

\newenvironment{keyidea}{\begin{block}{Key idea}}{\end{block}}
\newenvironment{activity}{\begin{block}{In class}}{\end{block}}
\newenvironment{cpdefinition}{\begin{block}{Definition}}{\end{block}}
\newenvironment{exampleblockcp}{
  \begingroup
  \setbeamercolor{block title}{bg=CalPolySand,fg=white}
  \setbeamercolor{block body}{bg=CalPolyBeige,fg=black}
  \begin{block}{Example}
}{
  \end{block}
  \endgroup
}

\newcommand{\imageplaceholder}[2][0.3in]{%
  \noindent\makebox[\linewidth][c]{%
    \fcolorbox{CalPolySage}{CalPolyBeige}{%
      \parbox[c][#1][c]{0.86\linewidth}{%
        \centering
        {\color{DeckMuted}\scriptsize Image placeholder: }%
        {\color{DeckInk}\scriptsize\ttfamily\detokenize{#2}}%
      }%
    }%
  }%
}

\newcommand{\sectionframe}[1]{
  \begingroup
  \setbeamercolor{background canvas}{bg=CalPolyGreen}
  \begin{frame}[plain]
    \vfill
    {\color{white}\LARGE\bfseries #1\par}
    \vspace{0.18in}
    {\color{CalPolyGold}\rule{0.22\paperwidth}{3pt}}
    {\color{white}\rule{0.62\paperwidth}{3pt}}
    \vfill
  \end{frame}
  \endgroup
}

\newcommand{\placeholdernote}{
  \begin{block}{Placeholder}
    Replace these bullets with the final lecture material, examples, and in-class activities.
  \end{block}
}


\pdfstringdefDisableCommands{\def\translate#1{#1}}

\graphicspath{{../assets/lecture-12/}}

% If an image file is not yet available, the slide displays a labeled
% placeholder using the intended filename. Place the matching file in
% ../assets/lecture-12/ to replace the placeholder automatically.
\newcommand{\lectureimage}[2][1.75in]{%
  \IfFileExists{../assets/lecture-12/#2}{%
    \includegraphics[width=\linewidth,height=#1,keepaspectratio]{#2}%
  }{%
    \fbox{%
      \parbox[c][#1][c]{0.94\linewidth}{%
        \centering
        \small\textbf{Image Placeholder}\\[0.08in]
        \scriptsize\texttt{\detokenize{#2}}
      }%
    }%
  }%
}

\setbeamersize{text margin left=0.42in,text margin right=0.42in}

\coursenumber{CS 1001}
\coursetitle{Introduction to Computing}
\lecturenumber{12}
\lecturetitle{Conditional Control Flow}
\lecturedate{September 21, 2026}
\courseunit{1. Computing and Program Execution}

\begin{document}

\makedecktitle

\begin{frame}{Today}
  \begin{itemize}
    \item Use Boolean expressions to choose whether statements execute.
    \item Trace \texttt{if}, \texttt{if}/\texttt{else}, and \texttt{if}/\texttt{elif}/\texttt{else} statements.
    \item Distinguish chained conditionals from nested conditionals.
    \item Use \texttt{input()} to obtain user data and convert it to an appropriate type.
    \item Apply Python's truth-value rules when conditions are not explicitly Boolean.
    \item Use a predicate function call directly as a condition.
  \end{itemize}

  \begin{keyidea}
    Conditional control flow changes \emph{which statement executes next} according to the truth value of an expression.
  \end{keyidea}
\end{frame}

\sectionframe{Part 1: From Boolean Values to Control Flow}

\begin{frame}{Sequential Execution Has Only One Path}
  \begin{columns}[T,onlytextwidth]
    \begin{column}{0.54\textwidth}
      \begin{block}{Sequential program}
        \texttt{x = 4}\\
        \texttt{y = x + 2}\\
        \texttt{z = y * 3}
      \end{block}

      \begin{itemize}
        \item Each statement executes after the preceding statement.
        \item Tracing asks two questions: what is the current state, and what executes next?
        \item Until now, the second question usually had one answer.
      \end{itemize}
    \end{column}

    \begin{column}{0.40\textwidth}
      % Intended visual: a straight control-flow arrow through three statements.
      \lectureimage[2.35in]{sequential-control-flow.png}
    \end{column}
  \end{columns}
\end{frame}

\begin{frame}{A Conditional Introduces a Choice of Execution Paths}
  \begin{cpdefinition}
    A \textbf{conditional statement} selects among possible execution paths according to the truth value of a condition.
  \end{cpdefinition}

  \begin{block}{Example}
    \texttt{temperature = 82}\\[0.06in]
    \texttt{if temperature > 80:}\\
    \hspace*{0.20in}\texttt{print("Warm day")}
  \end{block}

  \begin{itemize}
    \item \texttt{temperature > 80} is evaluated before the body is selected.
    \item The condition produces a truth value.
    \item That result determines whether the indented statement executes.
  \end{itemize}
\end{frame}

\begin{frame}{Reminder: Conditions Are Expressions}
  \begin{columns}[T,onlytextwidth]
    \begin{column}{0.47\textwidth}
      \begin{block}{Comparison}
        \texttt{age >= 18}\\[0.08in]
        $\longrightarrow$ \texttt{True} or \texttt{False}
      \end{block}

      \begin{block}{Compound Boolean expression}
        \texttt{age >= 18 and has\_id}
      \end{block}
    \end{column}

    \begin{column}{0.47\textwidth}
      \begin{itemize}
        \item Arithmetic may be evaluated inside a condition.
        \item Comparisons produce Boolean values.
        \item \texttt{and}, \texttt{or}, and \texttt{not} combine or negate Boolean claims.
        \item Parentheses may make grouping explicit.
      \end{itemize}
    \end{column}
  \end{columns}

  \begin{keyidea}
    The condition is evaluated using the same expression rules learned earlier; conditional control flow uses the resulting truth value.
  \end{keyidea}
\end{frame}

\begin{frame}{Reminder: Python Also Uses Truthiness}
  \begin{cpdefinition}
    \textbf{Truthiness} is Python's rule for interpreting a value in a Boolean context.
  \end{cpdefinition}

  \begin{center}
    \renewcommand{\arraystretch}{1.25}
    \begin{tabular}{l|l}
      \textbf{Common falsy values} & \textbf{Common truthy values} \\\hline
      \texttt{False} & \texttt{True} \\
      \texttt{0}, \texttt{0.0} & nonzero numbers \\
      \texttt{""} & nonempty strings \\
      \texttt{None} & most other ordinary values
    \end{tabular}
  \end{center}

  \begin{block}{Useful diagnostic}
    \centering
    \texttt{bool(value)} explicitly shows the truth value Python would use.
  \end{block}
\end{frame}

\begin{frame}{Tracing Now Requires State and Control}
  \begin{block}{Program}
    \texttt{x = 5}\\
    \texttt{if x > 3:}\\
    \hspace*{0.20in}\texttt{y = x + 1}\\
    \texttt{z = x * 2}
  \end{block}

  \begin{center}
    \renewcommand{\arraystretch}{1.20}
    \begin{tabular}{c|l|l}
      \textbf{Step} & \textbf{Executed item} & \textbf{Bindings afterward} \\\hline
      1 & \texttt{x = 5} & \texttt{x -> 5} \\
      2 & evaluate \texttt{x > 3} & condition is \texttt{True} \\
      3 & \texttt{y = x + 1} & \texttt{x -> 5, y -> 6} \\
      4 & \texttt{z = x * 2} & \texttt{x -> 5, y -> 6, z -> 10}
    \end{tabular}
  \end{center}

  \begin{keyidea}
    A trace must record not only values, but also which branch was selected.
  \end{keyidea}
\end{frame}

\sectionframe{Part 2: The \texttt{if} Statement}

\begin{frame}{The Simplest Conditional Is \texttt{if}}
  \begin{columns}[T,onlytextwidth]
    \begin{column}{0.54\textwidth}
      \begin{block}{General form}
        \texttt{if condition:}\\
        \hspace*{0.20in}\texttt{statement}\\
        \hspace*{0.20in}\texttt{statement}
      \end{block}

      \begin{itemize}
        \item Evaluate the condition.
        \item If its truth value is true, execute the indented suite.
        \item Otherwise, skip that suite.
        \item Continue with the first statement after the conditional.
      \end{itemize}
    \end{column}

    \begin{column}{0.40\textwidth}
      % Intended visual: decision diamond with true branch through body and false branch around it.
      \lectureimage[2.55in]{if-control-flow.png}
    \end{column}
  \end{columns}
\end{frame}

\begin{frame}{The Colon and Indentation Define the Conditional Suite}
  \begin{cpdefinition}
    A \textbf{suite} is the block of statements controlled by a compound statement such as \texttt{if}.
  \end{cpdefinition}

  \begin{block}{Example}
    \texttt{if score >= 70:}\\
    \hspace*{0.20in}\texttt{status = "passing"}\\
    \hspace*{0.20in}\texttt{message = "Keep going"}\\
    \texttt{print("Done")}
  \end{block}

  \begin{itemize}
    \item The colon introduces the suite.
    \item Both indented assignments belong to the \texttt{if}.
    \item \texttt{print("Done")} is outside the conditional and executes afterward regardless of the branch result.
  \end{itemize}
\end{frame}

\begin{frame}{A False Condition Skips the Entire Suite}
  \begin{block}{Program}
    \texttt{balance = 40}\\
    \texttt{if balance >= 100:}\\
    \hspace*{0.20in}\texttt{balance = balance - 100}\\
    \hspace*{0.20in}\texttt{print("Purchase complete")}\\
    \texttt{print(balance)}
  \end{block}

  \begin{enumerate}
    \item \texttt{balance >= 100} evaluates as \texttt{40 >= 100 -> False}.
    \item Neither indented statement executes.
    \item Control continues at \texttt{print(balance)}.
    \item The value of \texttt{balance} remains \texttt{40}.
  \end{enumerate}
\end{frame}

\begin{frame}{The Body May Contain Several Sequential Statements}
  \begin{block}{Program}
    \texttt{temperature = 95}\\
    \texttt{if temperature > 90:}\\
    \hspace*{0.20in}\texttt{message = "hot"}\\
    \hspace*{0.20in}\texttt{water = 2}\\
    \hspace*{0.20in}\texttt{print(message)}\\
    \texttt{print("Finished")}
  \end{block}

  \begin{itemize}
    \item The condition selects whether the \emph{entire} suite is entered.
    \item Once selected, statements inside the suite execute sequentially unless another control-flow construct changes that order.
    \item The first unindented statement marks the continuation point.
  \end{itemize}
\end{frame}

\begin{frame}{Trace Example: A Single \texttt{if}}
  \begin{block}{Program}
    \texttt{x = 6}\\
    \texttt{y = 2}\\
    \texttt{if x > y:}\\
    \hspace*{0.20in}\texttt{x = x - y}\\
    \hspace*{0.20in}\texttt{y = y + 1}\\
    \texttt{result = x + y}
  \end{block}

  \begin{center}
    \renewcommand{\arraystretch}{1.14}
    \begin{tabular}{c|l|l}
      \textbf{Step} & \textbf{Action} & \textbf{State} \\\hline
      1 & \texttt{x = 6} & \texttt{x -> 6} \\
      2 & \texttt{y = 2} & \texttt{x -> 6, y -> 2} \\
      3 & \texttt{x > y -> True} & branch selected \\
      4 & \texttt{x = x - y} & \texttt{x -> 4, y -> 2} \\
      5 & \texttt{y = y + 1} & \texttt{x -> 4, y -> 3} \\
      6 & \texttt{result = x + y} & \texttt{result -> 7}
    \end{tabular}
  \end{center}
\end{frame}

\begin{frame}{Prediction: Which Bindings Exist?}
  \begin{block}{Program}
    \texttt{points = 7}\\
    \texttt{if points >= 10:}\\
    \hspace*{0.20in}\texttt{bonus = 5}\\
    \hspace*{0.20in}\texttt{points = points + bonus}\\
    \texttt{final = points * 2}
  \end{block}

  \begin{activity}
    Before running Python, determine:
    \begin{enumerate}
      \item whether the conditional suite executes;
      \item whether \texttt{bonus} ever receives a binding;
      \item the final values of \texttt{points} and \texttt{final}.
    \end{enumerate}
  \end{activity}

  \begin{keyidea}
    A statement that is skipped cannot create or update a binding.
  \end{keyidea}
\end{frame}

\sectionframe{Part 3: Two-Way Selection with \texttt{if}/\texttt{else}}

\begin{frame}{\texttt{if}/\texttt{else} Selects Exactly One of Two Suites}
  \begin{columns}[T,onlytextwidth]
    \begin{column}{0.54\textwidth}
      \begin{block}{General form}
        \texttt{if condition:}\\
        \hspace*{0.20in}\texttt{true\_suite}\\
        \texttt{else:}\\
        \hspace*{0.20in}\texttt{false\_suite}
      \end{block}

      \begin{itemize}
        \item A true condition selects the first suite.
        \item A false condition selects the \texttt{else} suite.
        \item The two suites are mutually exclusive for one execution of the statement.
      \end{itemize}
    \end{column}

    \begin{column}{0.40\textwidth}
      % Intended visual: decision diamond with mutually exclusive true and false branches rejoining.
      \lectureimage[2.55in]{if-else-control-flow.png}
    \end{column}
  \end{columns}
\end{frame}

\begin{frame}{Both Branches Rejoin at the Continuation Point}
  \begin{block}{Program}
    \texttt{score = 84}\\
    \texttt{if score >= 70:}\\
    \hspace*{0.20in}\texttt{status = "pass"}\\
    \texttt{else:}\\
    \hspace*{0.20in}\texttt{status = "not pass"}\\
    \texttt{print(status)}
  \end{block}

  \begin{itemize}
    \item Exactly one assignment to \texttt{status} executes.
    \item After the selected branch finishes, both paths continue at \texttt{print(status)}.
    \item Designing both branches to establish the same required binding can simplify later code.
  \end{itemize}
\end{frame}

\begin{frame}{Trace Example: \texttt{if}/\texttt{else}}
  \begin{block}{Program}
    \texttt{x = 9}\\
    \texttt{if x \% 2 == 0:}\\
    \hspace*{0.20in}\texttt{kind = "even"}\\
    \texttt{else:}\\
    \hspace*{0.20in}\texttt{kind = "odd"}\\
    \texttt{length = len(kind)}
  \end{block}

  \begin{enumerate}
    \item \texttt{x \% 2 == 0 -> 1 == 0 -> False}.
    \item The first suite is skipped.
    \item The \texttt{else} suite binds \texttt{kind -> "odd"}.
    \item \texttt{len(kind) -> 3}, so \texttt{length -> 3}.
  \end{enumerate}
\end{frame}

\begin{frame}{Two Independent \texttt{if}s Are Not the Same as \texttt{if}/\texttt{else}}
  \begin{columns}[T,onlytextwidth]
    \begin{column}{0.47\textwidth}
      \begin{block}{Independent tests}
        \texttt{if x > 0:}\\
        \hspace*{0.16in}\texttt{print("positive")}\\
        \texttt{if x < 10:}\\
        \hspace*{0.16in}\texttt{print("small")}
      \end{block}

      Both tests are evaluated.
    \end{column}

    \begin{column}{0.47\textwidth}
      \begin{block}{Mutually exclusive alternatives}
        \texttt{if x > 0:}\\
        \hspace*{0.16in}\texttt{print("positive")}\\
        \texttt{else:}\\
        \hspace*{0.16in}\texttt{print("not positive")}
      \end{block}

      Exactly one suite is selected.
    \end{column}
  \end{columns}

  \begin{keyidea}
    Program structure expresses a logical relationship among alternatives; indentation is not merely visual formatting.
  \end{keyidea}
\end{frame}

\sectionframe{Part 4: Chained Conditionals with \texttt{elif}}

\begin{frame}{\texttt{elif} Represents Ordered Multi-Way Selection}
  \begin{columns}[T,onlytextwidth]
    \begin{column}{0.54\textwidth}
      \begin{block}{General form}
        \texttt{if condition1:}\\
        \hspace*{0.18in}\texttt{suite1}\\
        \texttt{elif condition2:}\\
        \hspace*{0.18in}\texttt{suite2}\\
        \texttt{elif condition3:}\\
        \hspace*{0.18in}\texttt{suite3}\\
        \texttt{else:}\\
        \hspace*{0.18in}\texttt{default\_suite}
      \end{block}
    \end{column}

    \begin{column}{0.40\textwidth}
      % Intended visual: ordered chain of decision diamonds, first true branch exits the chain.
      \lectureimage[2.65in]{elif-chain-control-flow.png}
    \end{column}
  \end{columns}
\end{frame}

\begin{frame}{Conditions in an \texttt{elif} Chain Are Tested Top to Bottom}
  \begin{enumerate}
    \item Evaluate the \texttt{if} condition.
    \item If it is truthy, execute that suite and skip the remainder of the chain.
    \item Otherwise, evaluate the first \texttt{elif} condition.
    \item Continue until one condition is truthy.
    \item If no preceding condition is truthy, execute \texttt{else} when present.
  \end{enumerate}

  \begin{keyidea}
    An \texttt{elif} chain selects the suite associated with the \textbf{first truthy condition}.
  \end{keyidea}
\end{frame}

\begin{frame}{Order Matters When Conditions Overlap}
  \begin{columns}[T,onlytextwidth]
    \begin{column}{0.47\textwidth}
      \begin{block}{Correctly ordered thresholds}
        \texttt{if score >= 90:}\\
        \hspace*{0.16in}\texttt{grade = "A"}\\
        \texttt{elif score >= 80:}\\
        \hspace*{0.16in}\texttt{grade = "B"}\\
        \texttt{elif score >= 70:}\\
        \hspace*{0.16in}\texttt{grade = "C"}
      \end{block}
    \end{column}

    \begin{column}{0.47\textwidth}
      \begin{alertblock}{Why order is semantic}
        A score of \texttt{95} satisfies all three comparisons. The first true condition determines the selected branch.
      \end{alertblock}

      For threshold classifications, testing from the most restrictive threshold downward is often the intended structure.
    \end{column}
  \end{columns}
\end{frame}

\begin{frame}{A Final \texttt{else} Can Make the Classification Exhaustive}
  \begin{block}{Example}
    \texttt{temperature = 58}\\[0.06in]
    \texttt{if temperature >= 90:}\\
    \hspace*{0.20in}\texttt{category = "hot"}\\
    \texttt{elif temperature >= 70:}\\
    \hspace*{0.20in}\texttt{category = "warm"}\\
    \texttt{elif temperature >= 50:}\\
    \hspace*{0.20in}\texttt{category = "cool"}\\
    \texttt{else:}\\
    \hspace*{0.20in}\texttt{category = "cold"}
  \end{block}

  \begin{itemize}
    \item Every numeric value reaches one branch.
    \item \texttt{58} fails the first two tests and satisfies \texttt{temperature >= 50}.
    \item The resulting binding is \texttt{category -> "cool"}.
  \end{itemize}
\end{frame}

\begin{frame}{Trace Example: A Chained Conditional}
  \begin{block}{Program}
    \texttt{x = 12}\\
    \texttt{if x < 0:}\\
    \hspace*{0.20in}\texttt{label = "negative"}\\
    \texttt{elif x == 0:}\\
    \hspace*{0.20in}\texttt{label = "zero"}\\
    \texttt{elif x < 10:}\\
    \hspace*{0.20in}\texttt{label = "small positive"}\\
    \texttt{else:}\\
    \hspace*{0.20in}\texttt{label = "large positive"}
  \end{block}

  \begin{block}{Evaluation sequence}
    \texttt{12 < 0 -> False} $\rightarrow$
    \texttt{12 == 0 -> False} $\rightarrow$
    \texttt{12 < 10 -> False} $\rightarrow$
    \texttt{else}
  \end{block}
\end{frame}

\begin{frame}{Boundary Values Are the Best Tests for Chained Conditions}
  \begin{activity}
    For the grading chain below, predict the branch selected for scores \texttt{69}, \texttt{70}, \texttt{79}, \texttt{80}, \texttt{89}, and \texttt{90}.
  \end{activity}

  \begin{block}{Chain}
    \texttt{if score >= 90:} \quad \texttt{grade = "A"}\\
    \texttt{elif score >= 80:} \quad \texttt{grade = "B"}\\
    \texttt{elif score >= 70:} \quad \texttt{grade = "C"}\\
    \texttt{else:} \quad \texttt{grade = "not passing"}
  \end{block}

  \begin{keyidea}
    Boundary cases expose mistakes involving \texttt{<} versus \texttt{<=}, ordering, and missing ranges.
  \end{keyidea}
\end{frame}

\sectionframe{Part 5: Nested Conditionals}

\begin{frame}{A Conditional Suite May Contain Another Conditional}
  \begin{columns}[T,onlytextwidth]
    \begin{column}{0.54\textwidth}
      \begin{block}{Example}
        \texttt{has\_ticket = True}\\
        \texttt{age = 15}\\[0.06in]
        \texttt{if has\_ticket:}\\
        \hspace*{0.18in}\texttt{if age >= 16:}\\
        \hspace*{0.36in}\texttt{status = "admit"}\\
        \hspace*{0.18in}\texttt{else:}\\
        \hspace*{0.36in}\texttt{status = "age restricted"}
      \end{block}
    \end{column}

    \begin{column}{0.40\textwidth}
      % Intended visual: one decision nested inside the true branch of another decision.
      \lectureimage[2.55in]{nested-conditional-tree.png}
    \end{column}
  \end{columns}
\end{frame}

\begin{frame}{Nested and Chained Conditionals Represent Different Structures}
  \begin{columns}[T,onlytextwidth]
    \begin{column}{0.47\textwidth}
      \begin{block}{Chained}
        \texttt{if A:}\\
        \hspace*{0.16in}\texttt{...}\\
        \texttt{elif B:}\\
        \hspace*{0.16in}\texttt{...}
      \end{block}

      \texttt{B} is considered only when \texttt{A} is false.
    \end{column}

    \begin{column}{0.47\textwidth}
      \begin{block}{Nested}
        \texttt{if A:}\\
        \hspace*{0.16in}\texttt{if B:}\\
        \hspace*{0.32in}\texttt{...}
      \end{block}

      \texttt{B} is considered only when \texttt{A} is true.
    \end{column}
  \end{columns}

  \begin{keyidea}
    Nesting expresses a dependency: reaching the inner decision depends on the outer path.
  \end{keyidea}
\end{frame}

\begin{frame}{Trace Nested Conditionals from the Outside In}
  \begin{block}{Program}
    \texttt{x = 7}\\
    \texttt{if x > 0:}\\
    \hspace*{0.20in}\texttt{sign = "positive"}\\
    \hspace*{0.20in}\texttt{if x \% 2 == 0:}\\
    \hspace*{0.40in}\texttt{kind = "even"}\\
    \hspace*{0.20in}\texttt{else:}\\
    \hspace*{0.40in}\texttt{kind = "odd"}\\
    \texttt{result = sign + " " + kind}
  \end{block}

  \begin{enumerate}
    \item Evaluate outer condition: \texttt{7 > 0 -> True}.
    \item Enter outer suite and bind \texttt{sign -> "positive"}.
    \item Evaluate inner condition: \texttt{7 \% 2 == 0 -> False}.
    \item Execute inner \texttt{else}: \texttt{kind -> "odd"}.
    \item Continue after the outer conditional: \texttt{result -> "positive odd"}.
  \end{enumerate}
\end{frame}

\begin{frame}{Choose the Structure That Matches the Logical Relationship}
  \begin{center}
    \renewcommand{\arraystretch}{1.30}
    \begin{tabular}{l|l}
      \textbf{Relationship among tests} & \textbf{Typical structure} \\\hline
      Perform something only when one claim is true & \texttt{if} \\
      Choose between two alternatives & \texttt{if}/\texttt{else} \\
      Choose the first matching case from an ordered set & \texttt{if}/\texttt{elif}/\texttt{else} \\
      A second decision is meaningful only within one outer case & nested \texttt{if}
    \end{tabular}
  \end{center}

  \begin{alertblock}{Design warning}
    Do not introduce nesting merely because it is available. Nest only when the logical dependency requires it.
  \end{alertblock}
\end{frame}

\sectionframe{Part 6: User Input, Truthiness, and Predicate Conditions}

\begin{frame}{\texttt{input()} Reads One Line of User Input as a String}
  \begin{columns}[T,onlytextwidth]
    \begin{column}{0.54\textwidth}
      \begin{cpdefinition}
        The built-in \texttt{input(prompt)} function displays the prompt, waits for a line of keyboard input, and returns the entered text as a \texttt{str}.
      \end{cpdefinition}

      \begin{block}{Example}
        \texttt{name = input("Name: ")}\\
        \texttt{print(name)}
      \end{block}

      Typing \texttt{Ada} produces the value \texttt{"Ada"}.
    \end{column}

    \begin{column}{0.40\textwidth}
      % Intended visual: keyboard input flowing through input() to a string value binding.
      \lectureimage[2.55in]{input-string-data-flow.png}
    \end{column}
  \end{columns}
\end{frame}

\begin{frame}{Numeric Input Must Be Converted Before Numeric Comparison}
  \begin{block}{Two-step form}
    \texttt{age\_text = input("Age: ")}\\
    \texttt{age = int(age\_text)}\\
    \texttt{if age >= 18:}\\
    \hspace*{0.20in}\texttt{print("Adult")}
  \end{block}

  \begin{itemize}
    \item If the user types \texttt{21}, \texttt{age\_text} is initially the string \texttt{"21"}.
    \item \texttt{int(age\_text)} evaluates to the integer \texttt{21} when the text is a valid integer representation.
    \item The comparison \texttt{age >= 18} is then a numeric comparison.
  \end{itemize}

  \begin{keyidea}
    User input is data; choose an appropriate representation before applying operations to it.
  \end{keyidea}
\end{frame}

\begin{frame}{Trace Example: Input Followed by a Conditional}
  \begin{block}{Program}
    \texttt{number\_text = input("Number: ")}\\
    \texttt{number = int(number\_text)}\\
    \texttt{if number > 0:}\\
    \hspace*{0.20in}\texttt{label = "positive"}\\
    \texttt{else:}\\
    \hspace*{0.20in}\texttt{label = "not positive"}\\
    \texttt{print(label)}
  \end{block}

  \begin{block}{Suppose the user enters \texttt{-3}}
    \texttt{number\_text -> "-3"} $\rightarrow$
    \texttt{number -> -3} $\rightarrow$
    \texttt{-3 > 0 -> False} $\rightarrow$
    \texttt{label -> "not positive"}
  \end{block}
\end{frame}

\begin{frame}{Truthiness Matters with User Input}
  \begin{block}{Program}
    \texttt{name = input("Name: ")}\\
    \texttt{if name:}\\
    \hspace*{0.20in}\texttt{message = "Hello, " + name}\\
    \texttt{else:}\\
    \hspace*{0.20in}\texttt{message = "No name entered"}
  \end{block}

  \begin{columns}[T,onlytextwidth]
    \begin{column}{0.47\textwidth}
      User enters \texttt{Ada}:\\
      \texttt{name -> "Ada"}\\
      nonempty string $\rightarrow$ truthy
    \end{column}

    \begin{column}{0.47\textwidth}
      User presses Enter immediately:\\
      \texttt{name -> ""}\\
      empty string $\rightarrow$ falsy
    \end{column}
  \end{columns}
\end{frame}

\begin{frame}[shrink=4]{A Predicate Function Call Can Be Used Directly as a Condition}
  \begin{cpdefinition}
    A \textbf{predicate function} is a function whose result represents a Boolean claim about its argument or arguments.
  \end{cpdefinition}

  \begin{columns}[T,onlytextwidth]
    \begin{column}{0.55\textwidth}
      \begin{block}{Example predicate}
        \texttt{def is\_even(n):}\\
        \hspace*{0.20in}\texttt{return n \% 2 == 0}
      \end{block}

      \begin{block}{Using the predicate}
        \texttt{value = 14}\\
        \texttt{if is\_even(value):}\\
        \hspace*{0.20in}\texttt{print("even")}
      \end{block}
    \end{column}

    \begin{column}{0.39\textwidth}
      % Intended visual: value -> predicate call -> True -> selected branch.
      \lectureimage[2.35in]{predicate-as-condition.png}
    \end{column}
  \end{columns}

  The call \texttt{is\_even(value)} is an expression whose returned Boolean value becomes the condition result.
\end{frame}

\begin{frame}[shrink=4]{Integrated Trace: Input, Chaining, and a Predicate}
  \begin{block}{Program}
    \texttt{def is\_even(n):}\\
    \hspace*{0.20in}\texttt{return n \% 2 == 0}\\[0.06in]
    \texttt{text = input("Integer: ")}\\
    \texttt{number = int(text)}\\
    \texttt{if number < 0:}\\
    \hspace*{0.20in}\texttt{category = "negative"}\\
    \texttt{elif number == 0:}\\
    \hspace*{0.20in}\texttt{category = "zero"}\\
    \texttt{elif is\_even(number):}\\
    \hspace*{0.20in}\texttt{category = "positive even"}\\
    \texttt{else:}\\
    \hspace*{0.20in}\texttt{category = "positive odd"}
  \end{block}

  \begin{activity}
    Trace the program for user inputs \texttt{-4}, \texttt{0}, \texttt{8}, and \texttt{7}. For each run, record the conditions evaluated, the first selected branch, and the final binding of \texttt{category}.
  \end{activity}

  \begin{keyidea}
    Trace conditionals by separating three questions: \textbf{what does the condition evaluate to, which path is selected, and how does that path change program state?}
  \end{keyidea}
\end{frame}

\end{document}
